% Options for packages loaded elsewhere
\PassOptionsToPackage{unicode}{hyperref}
\PassOptionsToPackage{hyphens}{url}
\documentclass[
  12pt,
  titlepage]{article}
\usepackage{xcolor}
\usepackage[a4paper, includeheadfoot, margin=2.54cm]{geometry}
\usepackage{amsmath,amssymb}
\setcounter{secnumdepth}{5}
\usepackage{iftex}
\ifPDFTeX
  \usepackage[T1]{fontenc}
  \usepackage[utf8]{inputenc}
  \usepackage{textcomp} % provide euro and other symbols
\else % if luatex or xetex
  \usepackage{unicode-math} % this also loads fontspec
  \defaultfontfeatures{Scale=MatchLowercase}
  \defaultfontfeatures[\rmfamily]{Ligatures=TeX,Scale=1}
\fi
\usepackage{lmodern}
\ifPDFTeX\else
  % xetex/luatex font selection
\fi
% Use upquote if available, for straight quotes in verbatim environments
\IfFileExists{upquote.sty}{\usepackage{upquote}}{}
\IfFileExists{microtype.sty}{% use microtype if available
  \usepackage[]{microtype}
  \UseMicrotypeSet[protrusion]{basicmath} % disable protrusion for tt fonts
}{}
\makeatletter
\@ifundefined{KOMAClassName}{% if non-KOMA class
  \IfFileExists{parskip.sty}{%
    \usepackage{parskip}
  }{% else
    \setlength{\parindent}{0pt}
    \setlength{\parskip}{6pt plus 2pt minus 1pt}}
}{% if KOMA class
  \KOMAoptions{parskip=half}}
\makeatother
\usepackage{longtable,booktabs,array}
\usepackage{caption}
\captionsetup[table]{skip=6pt}
\newcounter{none} % for unnumbered tables
\usepackage{calc} % for calculating minipage widths
% Correct order of tables after \paragraph or \subparagraph
\usepackage{etoolbox}
\makeatletter
\patchcmd\longtable{\par}{\if@noskipsec\mbox{}\fi\par}{}{}
\makeatother
% Allow footnotes in longtable head/foot
\IfFileExists{footnotehyper.sty}{\usepackage{footnotehyper}}{\usepackage{footnote}}
\makesavenoteenv{longtable}
\usepackage{graphicx}
\makeatletter
\newsavebox\pandoc@box
\newcommand*\pandocbounded[1]{% scales image to fit in text height/width
  \sbox\pandoc@box{#1}%
  \Gscale@div\@tempa{\textheight}{\dimexpr\ht\pandoc@box+\dp\pandoc@box\relax}%
  \Gscale@div\@tempb{\linewidth}{\wd\pandoc@box}%
  \ifdim\@tempb\p@<\@tempa\p@\let\@tempa\@tempb\fi% select the smaller of both
  \ifdim\@tempa\p@<\p@\scalebox{\@tempa}{\usebox\pandoc@box}%
  \else\usebox{\pandoc@box}%
  \fi%
}
% Set default figure placement to htbp
\def\fps@figure{htbp}
\makeatother
\setlength{\emergencystretch}{3em} % prevent overfull lines
\providecommand{\tightlist}{%
  \setlength{\itemsep}{0pt}\setlength{\parskip}{0pt}}
% ============================================================
% Preambulo moderno para Tarea Práctica 1
% Basado en la plantilla original del autor.
% Se conservan las interfaces: \materia, \TituloTarea, \Tony y \titulo.
% ============================================================

% Asegura que los .sty locales se encuentren dentro de la carpeta plantilla/
\makeatletter
%\def\input@path{{plantilla/}}
\makeatother

\usepackage{varios}
\graphicspath{{plantilla/}}
\usepackage{varios}
\graphicspath{{plantilla/}}
% ------------------------------------------------------------------
% Paleta Línea 12 (Línea Dorada) — sobrescribe los colores UNAM.
% Todo lo que ya usa estos nombres (portada, secciones, cajas,
% hipervínculos) hereda el dorado automáticamente.
% ------------------------------------------------------------------
\definecolor{UNAMBlue}{HTML}{B8901E}        % dorado principal
\definecolor{UNAMAccent}{HTML}{8A6A16}      % dorado oscuro (contraste)
\definecolor{UNAMLightBlue}{HTML}{F3E7C6}   % fondo claro dorado
\definecolor{UNAMLightAccent}{HTML}{EAD6A0} % fondo claro dorado oscuro
\definecolor{UNAMMuted}{HTML}{5B5B5B}       % gris neutro, sin cambio
\definecolor{UNAMText}{HTML}{262626}        % texto principal
\definecolor{UNAMLightGray}{HTML}{FBF9F4}   % blanco hueso

% Matemáticas modernas para XeLaTeX
\usepackage{unicode-math}
\setmathfont{TeX Gyre Termes Math}

%\spanishdecimal{.}

\renewcommand{\labelitemi}{$\bullet$}
\renewcommand{\labelitemii}{$\circ$}

% ============================================================
% Identidad del documento: conservar estos comandos
% ============================================================
\newcommand{\materia}{Modelos de Supervivencia y Series de Tiempo}
\newcommand{\TituloTarea}{Tarea Práctica 1}
\newcommand{\Tony}{\textsc{Anthonie Byron Ornelas Clement}}

% Autores
\newcommand{\Autores}{
    \textsc{Aguirre Garcia Abraham Isaias}\\
    \textsc{Larios Huitron Michelle}\\
    \textsc{Ornelas Clement Anthonie Byron}
}

% ============================================================
% Portada
% ============================================================



\newcommand{\titulo}{%
  \begin{minipage}{0.93\textwidth}
    \vspace{0.6cm}
    \begin{center}
      \IfFileExists{escudo.png}{%
        \includegraphics[height=2.7cm]{escudo.png}\par\vspace{0.55cm}
      }{%
        \IfFileExists{escudo0.png}{%
          \includegraphics[height=2.7cm]{escudo0.png}\par\vspace{0.55cm}
        }{}
      }
      {\sffamily\small\bfseries FACULTAD DE CIENCIAS\par}
      \vspace{0.12cm}
      {\sffamily\footnotesize\color{UNAMMuted}\materia\par}
      \vspace{0.75cm}
        \colorbox{UNAMBlue}{\parbox{0.82\textwidth}{%
        \centering
        \vspace{0.22cm}
        {\sffamily\color{white}\bfseries\fontsize{24}{29}\selectfont \TituloTarea\par}
        \vspace{0.22cm}
      }}
      \par\vspace{0.3cm}
      \textcolor{UNAMAccent}{\rule{4.5cm}{1pt}}\par
      \vspace{0.5cm}
      {\sffamily\normalsize\color{UNAMText}Análisis de la afluencia de pasajeros de la Línea 12 del Metro de la Ciudad de México\par}
      \vspace{0.85cm}
      {\sffamily\small\color{UNAMMuted}Período de análisis: operación plena, 2024--2025\par}
      \vfill
      %{\sffamily\small\bfseries \Tony\par}
      \vspace{0.85cm}
      {\large \Autores \par}
      \vspace{0.15cm}
    \end{center}
  \end{minipage}%
}

% ============================================================
% Ajustes visuales generales
% ============================================================
\setlength{\parindent}{0pt}
\setlength{\parskip}{0.55em}

% Hyperref lo administra Pandoc/R Markdown. Solo configuramos al inicio.
\makeatletter
\AtBeginDocument{%
  \@ifpackageloaded{hyperref}{%
    \hypersetup{
      colorlinks=true,
      linkcolor=UNAMBlue,
      urlcolor=UNAMBlue,
      citecolor=UNAMBlue
    }
  }{}
}
\makeatother

% Títulos
\usepackage{titlesec}
\titleformat{\section}
  {\sffamily\Large\bfseries\color{UNAMBlue}}
  {\thesection}{0.65em}{}
  [\vspace{0.18em}\color{UNAMAccent}\titlerule]

\titleformat{\subsection}
  {\sffamily\large\bfseries\color{UNAMText}}
  {\thesubsection}{0.6em}{}

\titleformat{\subsubsection}
  {\sffamily\normalsize\bfseries\color{UNAMBlue}}
  {\thesubsubsection}{0.55em}{}

% Figuras
\captionsetup{
  font=small,
  labelfont={bf,color=UNAMBlue},
  textfont={color=UNAMMuted},
  justification=centering,
  skip=7pt
}

% ============================================================
% Entornos de comunicación para el Rmd
% ============================================================
\newtcolorbox{queobservar}{%
  enhanced,
  breakable,
  colback=UNAMLightBlue,
  colframe=UNAMBlue,
  boxrule=0pt,
  leftrule=3pt,
  arc=2pt,
  left=8pt,right=8pt,top=7pt,bottom=7pt,
  before skip=8pt,after skip=10pt,
  fonttitle=\sffamily\bfseries,
  title=¿Qué debemos observar?%
}

\newtcolorbox{notatecnica}{%
  enhanced,
  breakable,
  colback=UNAMLightGray,
  colframe=UNAMMuted,
  boxrule=0pt,
  leftrule=3pt,
  arc=2pt,
  left=8pt,right=8pt,top=7pt,bottom=7pt,
  before skip=8pt,after skip=10pt,
  fonttitle=\sffamily\bfseries,
  title=Nota técnica%
}

\newtcolorbox{hallazgo}{%
  enhanced,
  breakable,
  colback=UNAMLightAccent,
  colframe=UNAMAccent,
  boxrule=0pt,
  leftrule=4pt,
  arc=2pt,
  left=8pt,right=8pt,top=8pt,bottom=8pt,
  before skip=9pt,after skip=11pt,
  fonttitle=\sffamily\bfseries,
  title=Hallazgo principal%
}


\newtcolorbox{citadestacada}{%
  enhanced,
  breakable,
  colback=white,
  colframe=UNAMBlue,
  boxrule=0pt,
  leftrule=2.4pt,
  arc=0pt,
  left=12pt,right=10pt,top=6pt,bottom=6pt,
  before skip=8pt,after skip=10pt,
  fontupper=\itshape\color{UNAMText}
}
\usepackage{booktabs}
\usepackage{longtable}
\usepackage{array}
\usepackage{multirow}
\usepackage{wrapfig}
\usepackage{float}
\usepackage{colortbl}
\usepackage{pdflscape}
\usepackage{tabularx}
\usepackage{xltabular}
\usepackage{threeparttable}
\usepackage{threeparttablex}
\usepackage[normalem]{ulem}
\usepackage{makecell}
\usepackage{xcolor}
\usepackage{bookmark}
\IfFileExists{xurl.sty}{\usepackage{xurl}}{} % add URL line breaks if available
\urlstyle{same}
% fallback for those not using the hyperref driver hyperxmp:
\makeatletter
\@ifundefined{xmpquote}{\newcommand{\xmpquote}[1]{#1}}{}
\makeatother
\hypersetup{
  hidelinks,
  pdfcreator={LaTeX via pandoc}}

\title{\titulo}
\author{}
\date{\vspace{-2.5em}2026-10-10}

\begin{document}
\maketitle

\section{Introducción}\label{introducciuxf3n}

\subsection{¿Qué queremos responder?}\label{quuxe9-queremos-responder}

La \textbf{Línea 12 del Sistema de Transporte Colectivo (STC) Metro de la Ciudad de México} conecta Mixcoac con Tláhuac a lo largo de 24.5 km y 20 estaciones. Por su importancia para la movilidad de la ciudad, cualquier trabajo de mantenimiento puede tener un efecto directo sobre un gran número de usuarios.

El propósito de este análisis no es solamente describir cuántas personas utilizan la línea. Buscamos identificar \textbf{cuándo la demanda es menor y qué patrones se repiten en el tiempo}, de modo que esa información pueda servir como apoyo para planear actividades de mantenimiento procurando reducir el impacto sobre los pasajeros.

La pregunta principal es:

\begin{citadestacada}
¿En qué momentos la Línea 12 presenta una menor afluencia de manera suficientemente consistente como para considerarlos una posible ventana de mantenimiento?
\end{citadestacada}

\subsection{Los datos}\label{los-datos}

Los datos corresponden a la \textbf{afluencia diaria por estación} publicados como datos abiertos por el Gobierno de la Ciudad de México. Para este trabajo se consideran los registros del \textbf{1 de enero de 2021 al 31 de diciembre de 2025}, para las 20 estaciones de la Línea 12.

Cada observación representa la cantidad de pasajeros registrada en una fecha. Para obtener una sola cifra diaria para toda la línea, se sumó la afluencia de las estaciones correspondientes a cada día.

Sin embargo, no sería adecuado tratar todo el período como si la línea hubiera operado bajo las mismas condiciones. El \textbf{3 de mayo de 2021} ocurrió el colapso de un tramo elevado entre las estaciones Olivos y Tezonco. Después del accidente, la línea permaneció cerrada y posteriormente reabrió por etapas:

\begingroup\fontsize{10}{12}\selectfont

\begin{longtable}[t]{|>{}c|>{}c|}
\hline
\cellcolor{UNAMBlue}{\textcolor{white}{\textbf{Fecha}}} & \cellcolor{UNAMBlue}{\textcolor{white}{\textbf{Estaciones.en.servicio}}}\\
\hline
\cellcolor{UNAMLightBlue}{15 ene 2023} & \cellcolor{UNAMLightBlue}{9 de 20}\\
\hline
\cellcolor{UNAMLightBlue}{16 jul 2023} & \cellcolor{UNAMLightBlue}{14 de 20}\\
\hline
\cellcolor{UNAMLightBlue}{23--29 ene 2024} & \cellcolor{UNAMLightBlue}{20 de 20}\\
\hline
\end{longtable}
\endgroup{}

Para estudiar el comportamiento habitual de la demanda, centraremos el análisis estadístico en el período de \textbf{operación plena}, del \textbf{30 de enero de 2024 al 31 de diciembre de 2025}. Este período contiene 702 días y corresponde a la etapa en la que las 20 estaciones estaban nuevamente en operación.

\subsection{¿Por qué usamos una serie de tiempo?}\label{por-quuxe9-usamos-una-serie-de-tiempo}

Una \textbf{serie de tiempo} es simplemente un conjunto de mediciones ordenadas por fecha. Su ventaja es que no solamente permite saber cuánto ocurrió, sino también observar \textbf{cómo cambia y qué comportamientos se repiten}.

En este caso nos interesan tres preguntas muy sencillas:

\begin{enumerate}
\def\labelenumi{\arabic{enumi}.}
\tightlist
\item
  ¿La cantidad de pasajeros está aumentando o disminuyendo con el tiempo?
\item
  ¿Hay días de la semana con una demanda sistemáticamente mayor o menor?
\item
  ¿Qué tan diferente es cada día respecto al comportamiento que normalmente esperaríamos?
\end{enumerate}

Estas preguntas guían los gráficos y transformaciones que se presentan en las siguientes secciones.

\section{Objetivo y metodología}\label{objetivo-y-metodologuxeda}

\subsection{Objetivo}\label{objetivo}

El objetivo es \textbf{identificar los principales patrones temporales de la afluencia de pasajeros de la Línea 12 durante su período de operación plena}, con énfasis en aquellos patrones que puedan ser útiles para la planeación de mantenimiento.

En particular, se busca:

\begin{itemize}
\tightlist
\item
  describir el comportamiento general de la afluencia;
\item
  identificar el patrón semanal de demanda;
\item
  separar la evolución en el tiempo del comportamiento que se repite cada semana; y
\item
  aplicar transformaciones que permitan estudiar los cambios de la serie de una forma más adecuada.
\end{itemize}

\subsection{Metodología}\label{metodologuxeda}

El procedimiento se puede resumir en cuatro pasos. Primero, se agregan los registros por estación para obtener la afluencia total de cada día. Después se selecciona el período de operación plena para evitar mezclar situaciones operativas muy distintas. A continuación, se utilizan gráficos y una descomposición de la serie para distinguir los patrones principales. Finalmente, se aplican transformaciones que ayudan a comparar los cambios de la demanda a lo largo del tiempo.

\begin{notatecnica}
El análisis de componentes se realiza mediante la descomposición clásica
\texttt{decompose()} y las transformaciones incluyen una escala logarítmica,
una diferencia de orden 1 y una diferencia estacional de periodo 7.
\end{notatecnica}

\section{Comportamiento de la afluencia}\label{comportamiento-de-la-afluencia}

\subsection{La historia de la Línea 12, 2021--2025}\label{la-historia-de-la-luxednea-12-20212025}

\begin{figure}

{\centering \includegraphics{Tarea_Practica_1_files/figure-latex/grafico-historia-1} 

}

\caption{Afluencia diaria de la Línea 12 entre 2021 y 2025. Las franjas distinguen las principales etapas de operación.}\label{fig:grafico-historia}
\end{figure}

\begin{queobservar}
La serie cambia de manera muy clara cuando cambia la situación operativa de la línea. Primero observamos la operación previa al accidente, después un cierre prolongado y finalmente una reapertura gradual hasta llegar a la operación con las 20 estaciones. Por esta razón, el análisis de patrones habituales se concentra en la etapa final.
\end{queobservar}

La figura también muestra que la afluencia durante 2024--2025 presenta una recuperación y una evolución ascendente. Para el problema planteado, esto es importante porque una misma diferencia absoluta de pasajeros puede representar una proporción distinta cuando el nivel general de demanda cambia.

\begin{notatecnica}
Las franjas del gráfico no forman parte del modelo; únicamente ayudan a contextualizar los cambios operativos que afectan a la interpretación de la serie.
\end{notatecnica}

\subsection{¿Qué día tiene menor demanda?}\label{quuxe9-duxeda-tiene-menor-demanda}

Una forma sencilla de comparar los días de la semana es calcular la \textbf{afluencia promedio para cada día} durante el período de operación plena. Esto permite responder directamente qué días concentran mayor o menor demanda.

\begin{figure}

{\centering \includegraphics{Tarea_Practica_1_files/figure-latex/grafico-semanal-1} 

}

\caption{Afluencia promedio por día de la semana durante la operación plena (2024--2025).}\label{fig:grafico-semanal}
\end{figure}

\begin{hallazgo}
El domingo destaca como la principal ventana de menor demanda. En promedio, la afluencia dominical es considerablemente menor que la observada en los días de mayor movimiento. El patrón semanal se repite a lo largo del período analizado, por lo que el domingo aparece como una alternativa natural para concentrar trabajos que reduzcan temporalmente la capacidad del servicio.
\end{hallazgo}

La utilidad para operaciones es sencilla: \textbf{si una intervención es inevitable, realizarla en un día de menor demanda potencialmente reduce el número de usuarios expuestos a la afectación}.

\begin{notatecnica}
El gráfico resume la serie mediante el promedio de afluencia para cada día de la semana. El análisis posterior permite estudiar este patrón semanal de manera conjunta con la tendencia de la serie.
\end{notatecnica}

\section{Hallazgos de la descomposición de la serie}\label{hallazgos-de-la-descomposiciuxf3n-de-la-serie}

\subsection{Separar lo que crece, lo que se repite y lo que cambia sin patrón claro}\label{separar-lo-que-crece-lo-que-se-repite-y-lo-que-cambia-sin-patruxf3n-claro}

Para interpretar la afluencia diaria, conviene separar tres tipos de comportamiento: lo que cambia de manera gradual con el tiempo, lo que se repite cada semana y las variaciones que no siguen un patrón evidente.

\begin{figure}

{\centering \includegraphics{Tarea_Practica_1_files/figure-latex/descomposicion-1} 

}

\caption{Descomposición clásica de la afluencia diaria durante la operación plena.}\label{fig:descomposicion}
\end{figure}

\textbf{Lo que crece con el tiempo.} La demanda aumenta durante el período de operación plena. Esto significa que la línea no está operando siempre al mismo nivel: con el paso del tiempo se incorporan más pasajeros al sistema.

\textbf{Lo que se repite cada semana.} El patrón semanal es estable y muestra claramente que los días de mayor demanda se concentran en la segunda mitad de la semana, mientras que el domingo presenta el nivel más bajo. Para la planeación operativa, este patrón repetitivo es especialmente valioso porque permite distinguir una baja esperable de una baja extraordinaria.

\textbf{Variaciones sin patrón claro.} Después de separar el crecimiento de la demanda y el ciclo semanal, las variaciones restantes representan movimientos que no siguen un patrón regular. Algunos días pueden alejarse del comportamiento esperado por vacaciones, días festivos, eventos u otras circunstancias específicas.

\begin{notatecnica}
Esta separación se obtuvo mediante la función
\texttt{decompose()} con una frecuencia de 7 y un modelo aditivo. El procedimiento separa la serie en tendencia, componente estacional semanal y variación residual.
\end{notatecnica}

\section{Hallazgos de las transformaciones aplicadas a la serie}\label{hallazgos-de-las-transformaciones-aplicadas-a-la-serie}

Las transformaciones no cambian la información original; \textbf{cambian la forma en que la observamos} para que sea más fácil estudiar ciertos aspectos de la serie.

\subsection{Poner los cambios en una escala comparable}\label{poner-los-cambios-en-una-escala-comparable}

Conforme aumenta la afluencia, también aumenta el tamaño de las diferencias absolutas entre días. Esto puede dificultar la comparación: un cambio de 30 mil pasajeros no tiene el mismo significado cuando la línea transporta 100 mil personas que cuando transporta 350 mil.

Por eso aplicamos una transformación logarítmica. En lugar de mirar únicamente cuántos pasajeros cambian, esta escala ayuda a interpretar los cambios \textbf{en relación con el nivel de afluencia}.

\begin{figure}

{\centering \includegraphics{Tarea_Practica_1_files/figure-latex/grafico-log-1} 

}

\caption{Afluencia diaria en su escala original (arriba) y en escala logarítmica (abajo).}\label{fig:grafico-log}
\end{figure}

\begin{hallazgo}
En la escala original, la amplitud entre un domingo de poca demanda y un viernes de mucha demanda crece junto con el nivel general de la línea: hacia el final del período esa diferencia es considerablemente mayor que al inicio. En la escala logarítmica, en cambio, el tamaño relativo de esa misma brecha semanal se mantiene mucho más parecido a lo largo del tiempo. En otras palabras: el domingo representa, proporcionalmente, un respiro de demanda similar en 2024 y en 2025, aunque en pasajeros absolutos la diferencia ya no sea la misma.
\end{hallazgo}

\begin{notatecnica}
La amplitud (máximo menos mínimo) calculada por bloques de 90 días pasa de alrededor de 225 mil pasajeros al inicio del período a más de 310 mil hacia el final, en la escala original. En la escala logarítmica, la misma amplitud se mantiene entre 1.0 y 1.5 en todos los bloques, una variación mucho menor en términos proporcionales. Esto es consistente con una varianza que crece junto con el nivel de la serie, el motivo habitual para aplicar una transformación logarítmica.
\end{notatecnica}

\subsection{¿Qué pasa al diferenciar la serie?}\label{quuxe9-pasa-al-diferenciar-la-serie}

Diferenciar significa comparar cada observación contra una observación anterior, en lugar de mirar el valor absoluto. Aquí usamos dos referencias distintas: el día inmediatamente anterior, y el mismo día de la semana una semana antes.

\begin{figure}

{\centering \includegraphics{Tarea_Practica_1_files/figure-latex/grafico-diff-1} 

}

\caption{Diferencia respecto al día anterior (arriba) y diferencia respecto al mismo día de la semana anterior (abajo).}\label{fig:grafico-diff}
\end{figure}

\begin{hallazgo}
Al comparar contra el día anterior (panel superior), el patrón semanal sigue siendo muy visible: los saltos de domingo a lunes y de viernes a sábado se repiten como picos y valles regulares, semana tras semana. Al comparar contra el mismo día de la semana anterior (panel inferior), en cambio, esos saltos regulares desaparecen casi por completo y la serie oscila de forma mucho más irregular alrededor de cero. Esto confirma que buena parte de lo que parece "ruido" en la serie original es, en realidad, el ciclo semanal predecible que ya identificamos: una vez que se compara cada día contra su propio equivalente de la semana pasada, lo que queda es mucho más parecido a variación sin patrón.
\end{hallazgo}

\begin{notatecnica}
La dispersión (desviación estándar) de la serie original es de aproximadamente 83 mil pasajeros. Al diferenciar contra el día anterior, la dispersión sube a cerca de 93 mil, porque se le suma el salto semanal a la variación normal. Al diferenciar contra el mismo día de la semana anterior, la dispersión baja a cerca de 48 mil: casi la mitad, señal de que esta transformación es la que mejor retira tanto la tendencia como la estacionalidad semanal, y la más adecuada como punto de partida para un modelo de pronóstico posterior.
\end{notatecnica}

\subsection{¿Qué aportan al análisis?}\label{quuxe9-aportan-al-anuxe1lisis}

En conjunto, estas transformaciones nos permiten mirar la misma información desde tres perspectivas: \textbf{comparar cambios proporcionales, detectar movimientos diarios y distinguir cambios inusuales de los que son parte del ciclo semanal}. Esto facilita la preparación de la serie para etapas posteriores de análisis y pronóstico.

\begin{notatecnica}
La primera transformación corresponde a $\log(y_t)$. La segunda corresponde a la primera diferencia, $y_t-y_{t-1}$, y la tercera a la diferencia estacional de periodo 7, $y_t-y_{t-7}$. Estas operaciones son habituales como preparación para modelos de series de tiempo porque ayudan a separar el comportamiento sistemático de los cambios restantes.
\end{notatecnica}

\section{Conclusiones}\label{conclusiones}

\begin{hallazgo}
El análisis permite llegar a cuatro conclusiones principales.

**1. La demanda cambia con el tiempo.** La operación plena de 2024--2025 presenta una evolución ascendente de la afluencia. Por ello, no es suficiente describir un nivel promedio: también importa considerar cómo cambia ese nivel a lo largo del tiempo.

**2. Existe un patrón semanal claro.** Los días de la semana no tienen la misma demanda. El domingo presenta la menor afluencia promedio y, en el extremo opuesto, los días de mayor actividad concentran una cantidad considerablemente mayor de pasajeros.

**3. El patrón semanal es útil para la planeación.** Debido a que las diferencias entre los días se repiten de forma consistente, el calendario semanal ofrece una señal sencilla para identificar posibles ventanas de menor demanda. Esto puede complementar la decisión operativa sobre cuándo realizar intervenciones de mantenimiento.

**4. Las transformaciones facilitan el análisis estadístico.** La escala logarítmica permite estudiar los cambios de forma proporcional, mientras que la diferenciación ayuda a retirar patrones persistentes para concentrarse en los cambios que no son explicados por el comportamiento habitual. En particular, la diferencia estacional de 7 días es la que deja la serie más cercana a un comportamiento sin patrón aparente, lo que la vuelve un punto de partida natural para etapas posteriores de modelación y pronóstico.
\end{hallazgo}

En términos prácticos, \textbf{el resultado más útil para la pregunta inicial es que el domingo aparece como una ventana recurrente de menor demanda}. No significa que toda intervención deba realizarse necesariamente ese día, pero sí que el comportamiento histórico de la serie proporciona evidencia para considerarlo dentro de la planeación, junto con otras restricciones operativas que no están incluidas en este análisis.

\begin{notatecnica}
Los resultados corresponden al período de operación plena definido en este informe. La etapa de modelación y pronóstico queda fuera del alcance de esta práctica, pero las transformaciones y hallazgos obtenidos constituyen una base para un modelo de series de tiempo posterior.
\end{notatecnica}

\section{Análisis de correlogramas (ACF y PACF)}\label{anuxe1lisis-de-correlogramas-acf-y-pacf}

\subsection{¿Para qué sirve un correlograma?}\label{para-quuxe9-sirve-un-correlograma}

Un \textbf{correlograma} muestra qué tan parecido es un valor de la serie a sus propios valores pasados, rezago por rezago (un día atrás, dos días atrás, y así sucesivamente). Hay dos versiones complementarias: la \textbf{ACF} (función de autocorrelación) mide esa relación de forma directa, y la \textbf{PACF} (autocorrelación parcial) mide lo mismo pero ``limpiando'' el efecto de los rezagos intermedios.

Esta herramienta es el paso natural antes de ajustar un modelo ARIMA: la forma de la ACF y la PACF sugiere qué tan complejo necesita ser el modelo y da una primera idea de sus parámetros, antes de afinarlos con una función de búsqueda automática como \texttt{auto.arima()}.

\subsection{Correlogramas de la serie original: la huella de la estacionalidad}\label{correlogramas-de-la-serie-original-la-huella-de-la-estacionalidad}

Antes de diferenciar, vale la pena ver cómo se ve el correlograma de la afluencia diaria tal cual, sin ningún ajuste.

\begin{figure}

{\centering \includegraphics{Tarea_Practica_1_files/figure-latex/correlograma-original-1} 

}

\caption{ACF y PACF de la afluencia diaria sin transformar, período de operación plena.}\label{fig:correlograma-original}
\end{figure}

\begin{hallazgo}
La ACF no decae: se mantiene alta incluso varias semanas atrás, y tiene picos muy marcados exactamente cada 7 rezagos. Esto es la firma típica de una serie con fuerte estacionalidad semanal todavía presente. Un correlograma con este aspecto no es útil para elegir directamente los parámetros de un modelo ARIMA, porque la estructura semanal domina por completo la lectura; primero hay que quitarla, como ya hicimos en la sección de transformaciones.
\end{hallazgo}

\begin{notatecnica}
En el rezago 1, ACF $\approx 0.37$ y PACF $\approx 0.37$. En los rezagos múltiplos de 7 (7, 14 y 21), la ACF alcanza $0.82$, $0.78$ y $0.75$ respectivamente, muy por encima de la banda de confianza aproximada de $\pm 1.96/\sqrt{n} \approx \pm 0.074$. Este patrón confirma, desde el correlograma, lo mismo que mostró la descomposición: la serie no es estacionaria y conserva un ciclo semanal muy fuerte.
\end{notatecnica}

\subsection{Correlogramas de la serie diferenciada: candidatos para el modelo ARIMA}\label{correlogramas-de-la-serie-diferenciada-candidatos-para-el-modelo-arima}

Al aplicar la diferencia estacional de 7 días que identificamos en la sección anterior, el correlograma cambia por completo.

\begin{figure}

{\centering \includegraphics{Tarea_Practica_1_files/figure-latex/correlograma-diff7-1} 

}

\caption{ACF y PACF de la serie diferenciada (lag 7), período de operación plena.}\label{fig:correlograma-diff7}
\end{figure}

\begin{hallazgo}
En los primeros rezagos (1 a 3), la ACF decae de forma gradual mientras que la PACF se corta de tajo después del rezago 1: esa combinación es la firma típica de un componente autorregresivo simple, un AR(1), para el comportamiento día a día. Además, tanto la ACF como la PACF muestran un quiebre marcado y negativo justo en el rezago 7 —consistente con haber diferenciado la serie con ese mismo periodo— lo que sugiere que conviene incluir también un componente de promedio móvil estacional. Fuera de esos puntos, el correlograma ya no muestra una estructura clara, señal de que la diferenciación sí dejó la serie en una forma mucho más manejable para un modelo ARIMA.
\end{hallazgo}

\begin{notatecnica}
En el rezago 1, ACF $\approx 0.37$ y PACF $\approx 0.37$; en los rezagos 2 y 3 la ACF cae a $0.17$ y $0.09$ mientras la PACF ya no es significativa ($\approx 0.04$ y $0.01$), un decaimiento geométrico consistente con un AR(1). En el rezago 7, ACF $\approx -0.37$ y PACF $\approx -0.38$, ambos muy por debajo de la banda de confianza. Como punto de partida razonable para explorar con \texttt{auto.arima()}, esto sugiere un modelo del tipo SARIMA(1,0,0)(0,1,1)$_{[7]}$: un término autorregresivo de orden 1 para la parte diaria, una diferencia estacional de periodo 7 (ya aplicada) y un término de promedio móvil estacional de orden 1. La forma más confiable de confirmar el orden final sigue siendo dejar que \texttt{auto.arima()} compare varias combinaciones por su criterio de información, usando este análisis como punto de partida informado en lugar de una búsqueda a ciegas.
\end{notatecnica}

Con este resultado, el equipo cuenta con una base concreta para la siguiente etapa del proyecto: ajustar y validar un modelo ARIMA sobre la serie de afluencia de la Línea 12.

\end{document}
